% ============================================================================
%  ESTADÍSTICA INFERENCIAL — Sesión 4
%  Semana 4: 31 de agosto al 5 de septiembre de 2026
%  1.5 Probabilidad condicional e independencia estadística
%  1.6 Fundamentos del muestreo
% ============================================================================
\documentclass[aspectratio=169,10pt]{beamer}

\usepackage[utf8]{inputenc}
\IfFileExists{spanish.ldf}{\usepackage[spanish,es-noshorthands]{babel}}{}
\usepackage{amsmath,amssymb}
\usepackage{booktabs}
\usepackage{tikz}
\usepackage{listings}

\usetheme{Uniminuto}

\uniasignatura{Estadística inferencial}
\uniprofesor{Andrés Felipe Puerta Velez}
\unicorreo{andres.puerta-v@uniminuto.edu.co}
\uninrc{CBASBL061}
\unisemana{Sesión 4 — Semana 4 · 31 de agosto al 5 de septiembre de 2026}

\newcommand{\listatiny}{\setbeamerfont{itemize/enumerate body}{size=\tiny}\tiny}

\newcommand{\fx}[1]{{\ttfamily\color{uniAzul}#1}}

% ------------------------------------------------------------- CÓDIGO R ----
\lstdefinestyle{Rcode}{
  language=R,
  basicstyle=\ttfamily\scriptsize,
  keywordstyle=\color{uniAzul}\bfseries,
  commentstyle=\color{uniGris}\itshape,
  stringstyle=\color{uniTitulo},
  showstringspaces=false, numbers=none, breaklines=true,
  columns=fullflexible,
  backgroundcolor=\color{uniCremaClaro},
  frame=leftline, framerule=1.2pt, rulecolor=\color{uniNaranja},
  xleftmargin=0.35cm, framexleftmargin=0.35cm,
  framexrightmargin=0.15cm, framextopmargin=2pt, framexbottommargin=2pt,
  aboveskip=0.15cm, belowskip=0.15cm
}
\lstdefinestyle{Rmini}{
  language=R,
  basicstyle=\ttfamily\tiny,
  keywordstyle=\color{uniAzul}\bfseries,
  commentstyle=\color{uniGris}\itshape,
  stringstyle=\color{uniTitulo},
  showstringspaces=false, numbers=none, breaklines=true,
  columns=fullflexible,
  backgroundcolor=\color{uniCremaClaro},
  frame=leftline, framerule=1.2pt, rulecolor=\color{uniNaranja},
  xleftmargin=0.35cm, framexleftmargin=0.35cm,
  framexrightmargin=0.15cm, framextopmargin=2pt, framexbottommargin=2pt,
  aboveskip=0.12cm, belowskip=0.12cm
}
\lstdefinestyle{Rout}{
  basicstyle=\ttfamily\tiny\color{black},
  showstringspaces=false, numbers=none, breaklines=true,
  columns=fullflexible,
  backgroundcolor=\color{uniGris!10},
  frame=leftline, framerule=1.2pt, rulecolor=\color{uniGris},
  xleftmargin=0.35cm, framexleftmargin=0.35cm,
  framexrightmargin=0.15cm, framextopmargin=2pt, framexbottommargin=2pt,
  aboveskip=0.12cm, belowskip=0.12cm
}

\begin{document}

\portada

\begin{frame}{Orden del día}
  \begin{enumerate}
    \item Llamado a lista y revisión de los ejercicios propuestos
    \item Repaso: dónde quedamos
    \item \azul{Tema 1.5: probabilidad condicional e independencia estadística}
    \item \azul{Teorema de Bayes: el caso de la prueba médica del 99\%}
    \item Actualización bayesiana: qué pasa al repetir la prueba
    \item \azul{Tema 1.6: fundamentos del muestreo}
    \item Ejercicios en clase
    \item Cierre y ejercicios propuestos
  \end{enumerate}
\end{frame}

\begin{frame}{Repaso: dónde quedamos}
  En la sesión 3 quedó pendiente una pieza de la regla de la suma:
  \[ P(A\cup B)=P(A)+P(B)-\azul{P(A\cap B)} \]

  \vspace{0.2cm}
  \begin{alertblock}{La pregunta de hoy}
    \footnotesize
    ¿Cómo se calcula $P(A\cap B)$ cuando no nos la dan? La respuesta es la
    \destacar{probabilidad condicional}: qué tanto cambia la probabilidad de un
    evento cuando ya sabemos que ocurrió otro.
  \end{alertblock}

  \vspace{0.2cm}
  {\footnotesize\color{uniGris} Esta es la herramienta que hace posible todo lo
  que viene después: si observar una muestra no cambiara nada sobre la
  población, la inferencia estadística no existiría.}
\end{frame}

\seccion{1.5 Probabilidad condicional e independencia}

\begin{frame}{Probabilidad condicional}
  \begin{block}{Definición}
    \[ P(A\mid B)=\frac{P(A\cap B)}{P(B)}, \qquad P(B)>0 \]
    Se lee ``probabilidad de $A$ \destacar{dado que} ocurrió $B$''.
  \end{block}

  \vspace{0.2cm}
  \begin{columns}[T,onlytextwidth]
    \begin{column}{0.46\textwidth}
      \centering
      \begin{tikzpicture}[scale=0.62]
        \draw[uniGris,thick] (-2.3,-1.5) rectangle (2.3,1.5);
        \fill[uniNaranja,opacity=0.22] (0.6,0) circle (1.15);
        \begin{scope}
          \clip (-0.6,0) circle (1.15);
          \fill[uniAzul,opacity=0.45] (0.6,0) circle (1.15);
        \end{scope}
        \draw[thick,uniAzul] (-0.6,0) circle (1.15);
        \draw[thick,uniNaranja] (0.6,0) circle (1.15);
        \node at (-1.35,0) {$A$}; \node at (1.35,0) {$B$};
        \node[font=\scriptsize] at (1.95,1.25) {$\Omega$};
      \end{tikzpicture}
    \end{column}
    \begin{column}{0.50\textwidth}
      \footnotesize
      Condicionar es \azul{cambiar el espacio muestral}: ya no miramos todo
      $\Omega$, sino solo $B$.

      \vspace{0.25cm}
      Por eso se divide entre $P(B)$: se está calculando qué fracción de $B$
      corresponde también a $A$.
    \end{column}
  \end{columns}
\end{frame}

\begin{frame}{Regla de la multiplicación}
  \begin{block}{Despejando de la definición}
    \[ P(A\cap B)=P(B)\,P(A\mid B)=P(A)\,P(B\mid A) \]
  \end{block}

  \vspace{0.2cm}
  \begin{exampleblock}{Ejemplo: extracción sin reemplazo}
    \footnotesize
    En una caja hay 12 productos, 4 defectuosos. Se extraen 2 sin devolver el
    primero:
    \[ P(\text{ambos defectuosos})=\frac{4}{12}\times\frac{3}{11}
       =\frac{12}{132}=0{,}0909 \]
    La segunda fracción ya está condicionada: quedan 11 productos y 3
    defectuosos.
  \end{exampleblock}

  \vspace{0.15cm}
  {\footnotesize\color{uniGris} Sin reemplazo las extracciones son
  dependientes; con reemplazo serían independientes.}
\end{frame}

\begin{frame}{Independencia estadística}
  \begin{block}{Definición}
    \small
    $A$ y $B$ son \destacar{independientes} si saber que ocurrió uno no cambia
    la probabilidad del otro:
    \[ P(A\mid B)=P(A) \iff P(A\cap B)=P(A)\,P(B) \]
  \end{block}

  \vspace{0.2cm}
  \begin{alertblock}{Independiente NO es lo mismo que excluyente}
    \footnotesize
    \begin{itemize}
      \item \azul{Excluyentes:} $P(A\cap B)=0$. Si ocurre uno, el otro es
            imposible. Son \emph{máximamente} dependientes.
      \item \azul{Independientes:} $P(A\cap B)=P(A)P(B)$. Ocurra o no uno, el
            otro no se entera.
    \end{itemize}
    Dos eventos con probabilidad positiva no pueden ser ambas cosas a la vez.
  \end{alertblock}
\end{frame}

\begin{frame}{Ejercicio 1 — en clase}
  \begin{alertblock}{En parejas (8 minutos)}
    De \destacar{500 empleados}, 200 están en Ventas y 300 en Operaciones. Han
    recibido capacitación 180: \destacar{90 de Ventas y 90 de Operaciones}. Se
    elige un empleado al azar.

    \vspace{0.25cm}
    \begin{enumerate}
      \item $P(V)$, $P(C)$ y $P(V\cap C)$
      \item $P(C\mid V)$ y $P(C\mid O)$
      \item $P(V\mid C)$
      \item ¿Son independientes ``estar en Ventas'' y ``estar capacitado''?
    \end{enumerate}
  \end{alertblock}
\end{frame}

\begin{frame}{Ejercicio 1 — solución}
  \begin{columns}[T,onlytextwidth]
    \begin{column}{0.42\textwidth}
      \footnotesize
      \begin{center}
      \begin{tabular}{@{}lccc@{}}
        \toprule
         & Capac. & No & Total \\
        \midrule
        Ventas      & 90  & 110 & 200 \\
        Operaciones & 90  & 210 & 300 \\
        \midrule
        Total       & 180 & 320 & 500 \\
        \bottomrule
      \end{tabular}
      \end{center}
      \vspace{0.2cm}
      {\scriptsize $P(V)=0{,}40$ \quad $P(C)=0{,}36$ \quad
       $P(V\cap C)=0{,}18$}
    \end{column}
    \begin{column}{0.54\textwidth}
      \begin{block}{Condicionales}
        \footnotesize
        \[ P(C\mid V)=\frac{90}{200}=\azul{0{,}450} \qquad
           P(C\mid O)=\frac{90}{300}=\azul{0{,}300} \]
        \[ P(V\mid C)=\frac{90}{180}=\azul{0{,}500} \]
      \end{block}
      \begin{alertblock}{Punto 4}
        \footnotesize
        $P(V)P(C)=0{,}40\times 0{,}36=0{,}144 \neq 0{,}18=P(V\cap C)$.
        \destacar{No son independientes:} en Ventas se capacita más.
      \end{alertblock}
    \end{column}
  \end{columns}
\end{frame}

\begin{frame}{Teorema de Bayes}
  \scriptsize
  \begin{block}{Para invertir el condicionamiento}
    \scriptsize
    \[ P(A\mid B)=\frac{P(B\mid A)\,P(A)}
                       {P(B\mid A)P(A)+P(B\mid A^{c})P(A^{c})} \]
    El denominador \azul{no es nuevo}: es la probabilidad total de $B$ que
    acabamos de ver. Bayes solo divide \destacar{el camino que nos interesa}
    entre \destacar{todos los caminos que llevan a $B$}.
  \end{block}

  \vspace{0.1cm}
  \begin{alertblock}{Por qué importa}
    \tiny
    Casi siempre se conoce $P(\text{síntoma}\mid\text{causa})$ —lo que mide un
    laboratorio— pero lo que se necesita para decidir es
    $P(\text{causa}\mid\text{síntoma})$. \destacar{No son lo mismo}, y
    confundirlos es el error que hoy vamos a desmontar con un caso real.
  \end{alertblock}

  \vspace{0.08cm}
  {\tiny\color{uniGris} El vocabulario que usaremos: $P(A)$ es la
  \azul{probabilidad a priori} (lo que se creía antes del dato), $P(B\mid A)$
  es la \azul{verosimilitud} (lo que aporta el dato) y $P(A\mid B)$ es la
  \azul{probabilidad a posteriori} (la creencia ya actualizada).}
\end{frame}

% ---------------------------------------------------------- CASO MÉDICO ----
\begin{frame}{El caso: una prueba médica ``99\% confiable''}
  \scriptsize
  \begin{block}{Los tres datos del problema}
    \tiny
    \begin{center}
    \begin{tabular}{@{}p{4.6cm}p{3.2cm}p{4.0cm}@{}}
      \toprule
      \textbf{Dato} & \textbf{En símbolos} & \textbf{Valor} \\
      \midrule
      La enfermedad afecta a 1 de cada 1\,000 personas &
        $P(D)$ & $0{,}001$ \ (0,1\%) \\
      Si usted está enfermo, la prueba acierta el 99\% &
        $P(+\mid D)$ & $0{,}99$ \ \azul{(sensibilidad)} \\
      Si usted está sano, la prueba se equivoca el 1\% &
        $P(+\mid D^{c})$ & $0{,}01$ \ \azul{(falso positivo)} \\
      \bottomrule
    \end{tabular}
    \end{center}
  \end{block}

  \vspace{0.1cm}
  \begin{alertblock}{La pregunta}
    \tiny
    Usted se hace la prueba \destacar{sin ningún síntoma}, por rutina, y sale
    \destacar{positiva}. ¿Cuál es la probabilidad de que realmente esté
    enfermo?

    \vspace{0.05cm}
    \azul{Antes de calcular, apueste un número.} Casi todo el mundo —incluidos
    médicos -- responde ``99\%''.
  \end{alertblock}
\end{frame}

\begin{frame}{Paso 1 — el árbol y la probabilidad total}
  \scriptsize
  \begin{columns}[T,onlytextwidth]
    \begin{column}{0.46\textwidth}
      \vspace{0.05cm}
      \begin{center}
      \begin{tikzpicture}[scale=0.80,every node/.style={font=\tiny}]
        \node (r) at (0,0) {};
        \node[uniAzul] (d)  at (2.0,1.5)  {$D$};
        \node[uniAzul] (nd) at (2.0,-1.5) {$D^{c}$};
        \node[uniTitulo] (dp) at (4.3,2.2)  {$+$};
        \node (dm) at (4.3,0.8)  {$-$};
        \node[uniTitulo] (np) at (4.3,-0.8) {$+$};
        \node (nm) at (4.3,-2.2) {$-$};
        \draw[uniGris] (r)--(d)  node[midway,above,sloped] {$0{,}001$};
        \draw[uniGris] (r)--(nd) node[midway,below,sloped] {$0{,}999$};
        \draw[uniGris] (d)--(dp) node[midway,above,sloped] {$0{,}99$};
        \draw[uniGris] (d)--(dm) node[midway,below,sloped] {$0{,}01$};
        \draw[uniGris] (nd)--(np) node[midway,above,sloped] {$0{,}01$};
        \draw[uniGris] (nd)--(nm) node[midway,below,sloped] {$0{,}99$};
        \node[uniTitulo,right] at (4.6,2.2)  {$0{,}00099$};
        \node[uniGris,right]   at (4.6,0.8)  {$0{,}00001$};
        \node[uniTitulo,right] at (4.6,-0.8) {$0{,}00999$};
        \node[uniGris,right]   at (4.6,-2.2) {$0{,}98901$};
      \end{tikzpicture}
      \end{center}
    \end{column}
    \begin{column}{0.51\textwidth}
      \begin{block}{Los dos caminos que dan positivo}
        \tiny
        \textbf{Enfermo y positivo:}
        \[ P(D)\cdot P(+\mid D) = 0{,}001\times 0{,}99
           = \azul{0{,}00099} \]
        \textbf{Sano y positivo (falsa alarma):}
        \[ P(D^{c})\cdot P(+\mid D^{c}) = 0{,}999\times 0{,}01
           = \azul{0{,}00999} \]
        \textbf{Probabilidad total de dar positivo:}
        \[ P(+) = 0{,}00099+0{,}00999 = \destacar{0{,}01098} \]
      \end{block}
    \end{column}
  \end{columns}

  \vspace{0.06cm}
  {\tiny\color{uniGris} \destacar{Fíjese en el desequilibrio:} las falsas
  alarmas ($0{,}00999$) son \azul{diez veces más} que las detecciones
  correctas ($0{,}00099$). No porque la prueba sea mala, sino porque
  \destacar{hay muchísima más gente sana que enferma}.}
\end{frame}

\begin{frame}{Paso 2 — Bayes, con cada número a la vista}
  \scriptsize
  \begin{exampleblock}{Sustituyendo en la fórmula}
    \tiny
    \[ P(D\mid +) \;=\;
       \frac{P(+\mid D)\,P(D)}
            {P(+\mid D)P(D)+P(+\mid D^{c})P(D^{c})} \]
    \[ =\;\frac{0{,}99\times 0{,}001}
              {\underbrace{0{,}99\times 0{,}001}_{0{,}00099}
               \;+\;\underbrace{0{,}01\times 0{,}999}_{0{,}00999}}
       \;=\;\frac{0{,}00099}{0{,}01098}
       \;=\;\destacar{0{,}0902} \]
  \end{exampleblock}

  \vspace{0.08cm}
  \begin{alertblock}{La respuesta: 9,02\%, no 99\%}
    \tiny
    Con una prueba que acierta el $99\%$ de las veces y un resultado
    positivo, la probabilidad de estar enfermo es de apenas un
    \destacar{9\%}. \azul{Más de nueve de cada diez positivos son falsas
    alarmas.}
  \end{alertblock}

  \vspace{0.06cm}
  {\tiny\color{uniGris} \destacar{De dónde viene la ilusión:} el $99\%$ es
  $P(+\mid D)$ —``si estoy enfermo, la prueba lo dice''—, pero lo que uno
  quiere saber es $P(D\mid +)$ —``la prueba lo dijo, ¿estoy enfermo?''—.
  \azul{Invertir la condición sin hacer la cuenta es la trampa.}}
\end{frame}

\begin{frame}{El mismo resultado, contando personas}
  \scriptsize
  \begin{block}{Imagine 100\,000 personas que se hacen la prueba}
    \tiny
    \begin{center}
    \begin{tabular}{@{}lrrr@{}}
      \toprule
      & \textbf{Prueba $+$} & \textbf{Prueba $-$} & \textbf{Total} \\
      \midrule
      Enfermos ($0{,}1\%$) &
        \azul{99} & 1 & 100 \\
      Sanos ($99{,}9\%$) &
        \azul{999} & 98\,901 & 99\,900 \\
      \midrule
      \textbf{Total} & \destacar{1\,098} & 98\,902 & 100\,000 \\
      \bottomrule
    \end{tabular}
    \end{center}
  \end{block}

  \vspace{0.1cm}
  \begin{exampleblock}{Leer la tabla es hacer Bayes}
    \tiny
    De las $1\,098$ personas que dieron positivo, solo $99$ están enfermas:
    \[ \frac{99}{1\,098} = \destacar{0{,}0902}
       \qquad\text{es decir, } \azul{1 de cada 11} \text{ positivos.} \]
  \end{exampleblock}

  \vspace{0.06cm}
  {\tiny\color{uniGris} \destacar{Esta tabla y la fórmula son lo mismo}, pero
  la tabla convence más rápido. Si alguna vez duda de un resultado bayesiano,
  \azul{invente una población grande y cuente cabezas}: es imposible
  equivocarse de intuición.}
\end{frame}

% ------------------------------------------------------ SEGUNDA PRUEBA ----
\begin{frame}{¿Y si repito la prueba? — el prior se actualiza}
  \scriptsize
  \begin{alertblock}{La idea clave de toda la inferencia}
    \tiny
    El médico pide una \destacar{segunda prueba en otro laboratorio} y vuelve
    a salir positiva. Ahora usted \azul{ya no es una persona cualquiera de la
    población}: es alguien que ya dio positivo una vez.

    \vspace{0.05cm}
    \destacar{El posterior de ayer es el prior de hoy:} se entra a la segunda
    cuenta con $P(D) = 0{,}0902$, no con $0{,}001$.
  \end{alertblock}

  \vspace{0.1cm}
  \begin{center}
  \begin{tikzpicture}[scale=0.95,every node/.style={font=\tiny,align=center},
    nodo/.style={draw=uniAzul,fill=uniAzul!10,rounded corners,inner sep=3pt,
                 minimum width=2.2cm,minimum height=0.6cm},
    flecha/.style={->,>=latex,uniGris,thick}]
    \node[nodo] (a) at (0,0) {Prior inicial\\ $0{,}001$};
    \node[nodo,fill=uniCrema,draw=uniTitulo] (b) at (3.4,0)
      {1.\textsuperscript{a} prueba $+$\\ $\longrightarrow 0{,}0902$};
    \node[nodo,fill=uniCrema,draw=uniTitulo] (c) at (7.0,0)
      {2.\textsuperscript{a} prueba $+$\\ $\longrightarrow 0{,}9075$};
    \node[nodo] (d) at (10.4,0) {Cada dato\\ mueve la creencia};
    \draw[flecha] (a)--(b); \draw[flecha] (b)--(c); \draw[flecha] (c)--(d);
  \end{tikzpicture}
  \end{center}

  \vspace{0.04cm}
  {\tiny\color{uniGris} Los datos \destacar{no se acumulan sumando}: cada uno
  se aplica sobre la creencia que dejó el anterior. Eso es
  \azul{actualización bayesiana}.}
\end{frame}

\begin{frame}{Segunda prueba — el cálculo completo}
  \scriptsize
  \begin{block}{Los datos de esta ronda}
    \tiny
    Prior: $P(D) = 0{,}0902$ \ (el posterior anterior). \quad
    Por lo tanto $P(D^{c}) = 1-0{,}0902 = 0{,}9098$. \quad
    La prueba no cambió: $P(+\mid D)=0{,}99$ y $P(+\mid D^{c})=0{,}01$.
  \end{block}

  \vspace{0.06cm}
  \begin{exampleblock}{Bayes otra vez, término por término}
    \tiny
    \textbf{Camino ``enfermo y positivo'':} \quad
    $0{,}0902\times 0{,}99 = \azul{0{,}08930}$

    \vspace{0.04cm}
    \textbf{Camino ``sano y positivo'':} \quad
    $0{,}9098\times 0{,}01 = \azul{0{,}00910}$

    \vspace{0.04cm}
    \textbf{Probabilidad total del segundo positivo:} \quad
    $0{,}08930+0{,}00910 = \destacar{0{,}09840}$
    \[ P(D\mid ++) \;=\; \frac{0{,}08930}{0{,}09840}
       \;=\; \destacar{0{,}9075} \]
  \end{exampleblock}

  \vspace{0.06cm}
  {\tiny\color{uniGris} \destacar{De 9\% a 90,75\% con el mismo dato
  repetido.} La segunda prueba no es más precisa que la primera: lo que cambió
  es que \azul{ya no se compara contra 999 sanos, sino contra 9}. El
  desequilibrio que arruinaba la primera cuenta casi desapareció.}
\end{frame}

\begin{frame}{Dos comprobaciones que vale la pena hacer}
  \scriptsize
  \begin{columns}[T,onlytextwidth]
    \begin{column}{0.49\textwidth}
      \begin{exampleblock}{1. Las dos rutas coinciden}
        \tiny
        Si se tratan las dos pruebas \azul{de una sola vez}, con
        $P(++\mid D) = 0{,}99^{2}$ y $P(++\mid D^{c}) = 0{,}01^{2}$:
        \[ \frac{0{,}001\times 0{,}9801}
                {0{,}001\times 0{,}9801 + 0{,}999\times 0{,}0001} \]
        \[ = \frac{0{,}0009801}{0{,}00108}
           = \destacar{0{,}9075} \ \checkmark \]
        {\color{uniGris}Idéntico. Actualizar por partes o de golpe da lo
        mismo.}
      \end{exampleblock}
    \end{column}
    \begin{column}{0.49\textwidth}
      \begin{alertblock}{2. ¿Y si la segunda da negativo?}
        \tiny
        Prior $0{,}0902$; ahora interesa $P(-\mid D)=0{,}01$ y
        $P(-\mid D^{c})=0{,}99$:
        \[ \frac{0{,}0902\times 0{,}01}
                {0{,}0902\times 0{,}01+0{,}9098\times 0{,}99} \]
        \[ = \frac{0{,}000902}{0{,}901604}
           = \destacar{0{,}001} \]
        {\color{uniGris}Vuelve \azul{exactamente} al 0,1\% de partida: un
        positivo y un negativo se cancelan.}
      \end{alertblock}
    \end{column}
  \end{columns}

  \vspace{0.08cm}
  \begin{block}{La moraleja para el resto del curso}
    \tiny
    \destacar{Un dato aislado casi nunca decide; lo que decide es la
    acumulación de datos sobre una creencia previa honesta.} Por eso ningún
    diagnóstico serio se hace con una sola prueba, y por eso en la unidad 3 no
    aceptaremos una conclusión por un solo valor $p$ sin mirar el contexto.
  \end{block}
\end{frame}

\begin{frame}{Ejercicio 2 — en clase}
  \begin{alertblock}{En parejas (12 minutos). Escriba cada producto.}
    \scriptsize
    Una prueba de tamizaje para una condición que afecta al \destacar{2\%}
    de la población detecta al \destacar{95\%} de quienes la tienen, pero
    marca como positivo al \destacar{8\%} de quienes no la tienen.

    \vspace{0.2cm}
    \begin{enumerate}\setlength{\itemsep}{0.02cm}
      \item ¿Cuál es la probabilidad de que una persona cualquiera dé
            positivo?
      \item Si una persona da positivo, ¿cuál es la probabilidad de que
            realmente tenga la condición?
      \item Le repiten la prueba y vuelve a dar positivo.
            \azul{Actualice} la probabilidad.
      \item Interprete los dos resultados en una frase.
    \end{enumerate}
  \end{alertblock}
\end{frame}

\begin{frame}{Ejercicio 2 — solución}
  \scriptsize
  \begin{columns}[T,onlytextwidth]
    \begin{column}{0.49\textwidth}
      \begin{block}{Puntos 1 y 2 — primera prueba}
        \tiny
        \textbf{1.} $P(+) = \underbrace{0{,}02\times 0{,}95}_{0{,}0190}
        + \underbrace{0{,}98\times 0{,}08}_{0{,}0784}
        = \azul{0{,}0974}$

        \vspace{0.1cm}
        \textbf{2.} $P(C\mid +) = \dfrac{0{,}0190}{0{,}0974}
        = \destacar{0{,}1951}$ \ (19,51\%)
      \end{block}
    \end{column}
    \begin{column}{0.49\textwidth}
      \begin{block}{Punto 3 — segunda prueba}
        \tiny
        Prior nuevo: $0{,}1951$; \ $P(C^{c}) = 0{,}8049$.

        \vspace{0.06cm}
        $0{,}1951\times 0{,}95 = \azul{0{,}18535}$

        $0{,}8049\times 0{,}08 = \azul{0{,}06439}$

        Total $= \destacar{0{,}24974}$
        \[ P(C\mid ++) = \frac{0{,}18535}{0{,}24974}
           \approx \destacar{0{,}742} \]
      \end{block}
    \end{column}
  \end{columns}

  \vspace{0.08cm}
  \begin{alertblock}{Punto 4 — interpretación}
    \tiny
    Con un positivo, \destacar{cuatro de cada cinco} personas marcadas están
    sanas: no es motivo para alarmarse, es motivo para \azul{confirmar}. Con
    el segundo positivo la probabilidad sube al $74\%$ y ya sí cambia la
    conducta clínica. \destacar{El tamizaje no diagnostica: solo decide a
    quién vale la pena estudiar mejor.}
  \end{alertblock}
\end{frame}

\seccion{1.6 Fundamentos del muestreo}

\begin{frame}{Conceptos básicos del muestreo}
  \footnotesize
  \begin{center}
  \begin{tabular}{@{}p{2.9cm}p{8.4cm}@{}}
    \toprule
    \textbf{Término} & \textbf{Significado} \\
    \midrule
    Población & Conjunto completo de elementos sobre los que se quiere
      concluir. Su tamaño se denota $N$. \\
    Muestra & Subconjunto de la población que sí se observa. Su tamaño es $n$. \\
    Unidad de muestreo & El elemento individual que se selecciona: un cliente,
      un pedido, una factura. \\
    Marco muestral & La lista de la que se extrae la muestra: base de datos,
      censo, registro. \\
    Parámetro & Número que describe a la \azul{población}: $\mu$, $\sigma$, $p$.
      Es fijo y desconocido. \\
    Estadístico & Número calculado en la \azul{muestra}: $\bar{x}$, $s$,
      $\hat{p}$. Cambia con cada muestra. \\
    Error muestral & Diferencia entre el estadístico y el parámetro, debida a
      que se observó solo una parte. \\
    \bottomrule
  \end{tabular}
  \end{center}
\end{frame}

\begin{frame}{Parámetro y estadístico: la distinción central del curso}
  \begin{columns}[T,onlytextwidth]
    \begin{column}{0.48\textwidth}
      \begin{block}{Población ($N$)}
        \footnotesize
        \begin{itemize}
          \item Media $\mu$
          \item Desviación $\sigma$
          \item Proporción $p$
        \end{itemize}
        Valores \azul{fijos} que casi nunca conocemos.
      \end{block}
    \end{column}
    \begin{column}{0.48\textwidth}
      \begin{exampleblock}{Muestra ($n$)}
        \footnotesize
        \begin{itemize}
          \item Media $\bar{x}$
          \item Desviación $s$
          \item Proporción $\hat{p}$
        \end{itemize}
        Valores \azul{aleatorios}: cambian de muestra en muestra.
      \end{exampleblock}
    \end{column}
  \end{columns}

  \vspace{0.3cm}
  \begin{alertblock}{Todo el curso en una frase}
    \footnotesize
    Usamos un estadístico \emph{aleatorio} para decir algo sobre un parámetro
    \emph{fijo} y desconocido, y cuantificamos con probabilidad qué tan lejos
    podemos estar. \destacar{Esa cuantificación es la inferencia estadística.}
  \end{alertblock}
\end{frame}

\begin{frame}{Por qué muestrear en vez de censar}
  \begin{columns}[T,onlytextwidth]
    \begin{column}{0.48\textwidth}
      \begin{block}{Razones}
        \footnotesize
        \begin{itemize}
          \item Costo y tiempo
          \item Poblaciones muy grandes o infinitas
          \item Pruebas destructivas (no se puede quemar todo el lote)
          \item Con buen diseño, la precisión es suficiente
        \end{itemize}
      \end{block}
    \end{column}
    \begin{column}{0.48\textwidth}
      \begin{alertblock}{Riesgos}
        \footnotesize
        \begin{itemize}
          \item Marco muestral incompleto
          \item Sesgo de selección
          \item No respuesta
          \item Muestra demasiado pequeña
        \end{itemize}
      \end{alertblock}
    \end{column}
  \end{columns}

  \vspace{0.25cm}
  {\footnotesize\color{uniGris} Una muestra grande mal elegida es peor que una
  pequeña bien elegida: el tamaño reduce el error \emph{aleatorio}, pero no
  corrige el \emph{sesgo}. En la sesión 5 veremos cómo elegir bien.}
\end{frame}

\begin{frame}{Ejercicio 3 — en clase}
  \begin{alertblock}{Individual (7 minutos)}
    Una caja contiene \destacar{10 bolas: 6 blancas y 4 negras}. Se extraen 2
    al azar \emph{sin reemplazo}.

    \vspace{0.25cm}
    \begin{enumerate}
      \item $P(\text{ambas blancas})$
      \item $P(\text{ambas negras})$
      \item $P(\text{una de cada color})$
      \item $P(\text{al menos una blanca})$
      \item Si la primera fue blanca, ¿cuál es la probabilidad de que la
            segunda también lo sea?
    \end{enumerate}
  \end{alertblock}
\end{frame}

\begin{frame}{Ejercicio 3 — solución}
  \begin{columns}[T,onlytextwidth]
    \begin{column}{0.52\textwidth}
      \begin{block}{Cálculos}
        \footnotesize
        \begin{align*}
          P(BB) &= \tfrac{6}{10}\cdot\tfrac{5}{9}=\azul{0{,}3333}\\[0.2em]
          P(NN) &= \tfrac{4}{10}\cdot\tfrac{3}{9}=\azul{0{,}1333}\\[0.2em]
          P(\text{una de cada}) &= \tfrac{6}{10}\tfrac{4}{9}+\tfrac{4}{10}\tfrac{6}{9}
             =\azul{0{,}5333}\\[0.2em]
          P(\ge 1\,\text{blanca}) &= 1-P(NN)=\azul{0{,}8667}\\[0.2em]
          P(B_2\mid B_1) &= \tfrac{5}{9}=\azul{0{,}5556}
        \end{align*}
      \end{block}
    \end{column}
    \begin{column}{0.44\textwidth}
      \begin{exampleblock}{Verificación}
        \footnotesize
        Los tres primeros casos son excluyentes y agotan todo:
        \[ 0{,}3333+0{,}1333+0{,}5333=1 \]
      \end{exampleblock}
      \begin{alertblock}{Punto 5}
        \footnotesize
        $0{,}5556 \neq 0{,}60=P(B_1)$: sacar una blanca \azul{sí} cambia la
        probabilidad de la siguiente. Sin reemplazo no hay independencia.
      \end{alertblock}
    \end{column}
  \end{columns}
\end{frame}

\begin{frame}[fragile]{Probabilidad condicional y muestreo, en R}
  \scriptsize
\begin{lstlisting}[style=Rmini]
# Una tabla de doble entrada y sus probabilidades
tabla <- matrix(c(30, 20, 25, 25), nrow = 2,
                dimnames = list(Genero = c("M", "F"),
                                Compra = c("Si", "No")))
prop.table(tabla)                  # probabilidades conjuntas
prop.table(tabla, margin = 1)      # condicionales P(Compra | Genero)
addmargins(tabla)                  # con totales

# Muestreo aleatorio simple
set.seed(2026)                     # para que el resultado sea reproducible
sample(1:500, size = 30)           # 30 unidades sin reemplazo
\end{lstlisting}
\begin{lstlisting}[style=Rout]
      Compra
Genero  Si  No
     M 0.6 0.4
     F 0.5 0.5
\end{lstlisting}

  \vspace{-0.1cm}
  \begin{alertblock}{Dos detalles}
    \tiny
    \fx{margin = 1} condiciona \destacar{por filas}: cada fila suma 1 y es
    justamente $P(\text{Compra} \mid \text{Género})$. \quad
    \fx{set.seed()} fija el azar: sin él, cada corrida da una muestra distinta
    y el trabajo deja de ser reproducible.
  \end{alertblock}
\end{frame}

\begin{frame}{Cierre de la sesión}
  \begin{block}{Lo que debe quedar claro hoy}
    \begin{itemize}
      \item $P(A\mid B)=P(A\cap B)/P(B)$: condicionar es reducir el espacio
            muestral a $B$.
      \item $P(A\cap B)=P(B)P(A\mid B)$, y si son independientes
            $P(A\cap B)=P(A)P(B)$.
      \item Independiente $\neq$ excluyente.
      \item Parámetro (población, fijo) frente a estadístico (muestra,
            aleatorio).
    \end{itemize}
  \end{block}

  \vspace{0.2cm}
  \begin{alertblock}{Trabajo independiente para la próxima sesión}
    \footnotesize
    Resolver los \destacar{20 ejercicios de la lámina siguiente}; las
    respuestas están en la lámina posterior. La sesión 5 cierra la unidad 1 con
    los tipos de muestreo y el tamaño de la muestra, y hay repaso para el
    examen 1.
  \end{alertblock}
\end{frame}

\begin{frame}{Ejercicios propuestos}
  \vspace{-0.1cm}
  \begin{beamercolorbox}[wd=\textwidth,sep=0.14cm]{block body}
    \centering\scriptsize
    \textbf{Del 1 al 10:} 600 clientes. Canal: 360 usan la app ($A$), 240 van a
    la oficina. Producto: 210 tienen crédito ($C$), de los cuales 150 usan la app.
  \end{beamercolorbox}
  \vspace{0.05cm}
  \listatiny
  \vspace{0.15cm}
  \begin{columns}[T,onlytextwidth]
    \begin{column}{0.47\textwidth}
      \begin{enumerate}\setlength{\itemsep}{0.16cm}
        \item $P(A)$
        \item $P(C)$
        \item $P(A\cap C)$
        \item $P(C\mid A)$
        \item $P(C\mid A^{c})$
        \item $P(A\mid C)$
        \item $P(A\cap C^{c})$
        \item $P(C^{c})$
        \item $P(A)\cdot P(C)$
        \item ¿Son $A$ y $C$ independientes?
      \end{enumerate}
    \end{column}
    \begin{column}{0.50\textwidth}
      \begin{enumerate}\setlength{\itemsep}{0.16cm}
        \setcounter{enumi}{10}
        \item Urna (6 blancas, 4 negras, sin reemplazo): $P(\text{2 blancas})$
        \item $P(\text{2 negras})$
        \item $P(\text{una de cada color})$
        \item $P(\text{al menos una blanca})$
        \item $P(B_2\mid B_1)$
        \item Dos dados: $P(\text{ambos } 6)$
        \item Dos dados: $P(\text{2.º es }6 \mid \text{1.º es }6)$
        \item Dos monedas: $P(\text{ambas cara})$
        \item Prueba: 2\% defectuosas, 95\% detección, 4\% falsos
              positivos. $P(T^{+})$
        \item Del 19: $P(D\mid T^{+})$
      \end{enumerate}
    \end{column}
  \end{columns}
\end{frame}

\begin{frame}{Respuestas}{Autoevaluación de los ejercicios propuestos}
  \listatiny
  \vspace{0.15cm}
  \begin{columns}[T,onlytextwidth]
    \begin{column}{0.47\textwidth}
      \begin{enumerate}\setlength{\itemsep}{0.16cm}
        \item $0{,}6000$
        \item $0{,}3500$
        \item $0{,}2500$
        \item $150/360=0{,}4167$
        \item $60/240=0{,}2500$
        \item $150/210=0{,}7143$
        \item $210/600=0{,}3500$
        \item $0{,}6500$
        \item $0{,}2100$
        \item No: $0{,}2100\neq 0{,}2500$
      \end{enumerate}
    \end{column}
    \begin{column}{0.50\textwidth}
      \begin{enumerate}\setlength{\itemsep}{0.16cm}
        \setcounter{enumi}{10}
        \item $0{,}3333$
        \item $0{,}1333$
        \item $0{,}5333$
        \item $0{,}8667$
        \item $0{,}5556$
        \item $0{,}0278$
        \item $0{,}1667$ \;{\color{uniTitulo}independientes}
        \item $0{,}2500$
        \item $0{,}0582$
        \item $0{,}3265$
      \end{enumerate}
    \end{column}
  \end{columns}
\end{frame}

\begin{frame}{Referencias}
  \footnotesize
  \begin{itemize}
    \item[{[1]}] Díaz Rodríguez, M. (2022). \textit{Estadística inferencial
      aplicada}. Universidad del Norte.
      \texttt{https://elibro.net/es/lc/uniminuto/titulos/221614}
    \item[{[2]}] Llinás Solano, H. (2014). \textit{Introducción a la
      estadística con aplicaciones en ciencias sociales}. Universidad del
      Norte. \texttt{https://elibro.net/es/lc/uniminuto/titulos/70062}
    \item[{[3]}] Contento Rubio, M. R. (2019). \textit{Estadística con
      aplicaciones en R}. Editorial Utadeo.
      \texttt{https://elibro.net/es/lc/uniminuto/titulos/220926}
  \end{itemize}

  \vspace{0.35cm}
  \begin{center}
    {\Large\color{uniAzul}\textbf{¿Preguntas?}}
  \end{center}
\end{frame}

\end{document}
